\chapter{What Happens Between Measurements}
\label{chap:what-happens-between-measurements}

% Closed question: What can be said about the gap between recorded events
% without inventing structure not in the ledger?

\begin{chapteressay}

A highway camera records at ten frames per second. A car enters the frame in
one image, moves across the lane, and appears again one tenth of a second later.
The camera has two records. It does not have the car between them.

The missing interval is not empty. The car did not vanish and reappear. It
traveled through the gap. But the camera did not record each intermediate
position. If an investigator wants to know whether the car crossed the lane
marker, the two frames constrain the answer, but they do not contain the whole
motion. The investigator may use the frame rate, visible positions, motion blur,
known vehicle dynamics, and road geometry. The investigator may not simply
invent a path because it would be convenient.

That is the structure of this chapter: the gap is real, and the gap is
constrained.

The speedometer has the same problem. At 60 mph, the Hall-effect sensor may
admit hundreds of pulses per second, but still not infinitely many. Between two
adjacent tooth crossings, the instrument has no direct mark. If the driver taps
the brake and accelerates again inside that interval, the speedometer may not
carry the event. The displayed speed is an interpolation over admitted marks,
not a perfect film of motion.

The same structure appears in medicine. A patient wears a glucose monitor that
reports a value every five minutes. Between readings, the patient's blood sugar
changes. A meal, a walk, or insulin may bend the curve inside the interval. The
monitor's graph is not the body itself. It is a record at a sampling rate. The
curve drawn between points is a constrained continuation, not a new
measurement. It may be useful; it may even be the best admissible estimate. But
it must not pretend to contain events that the instrument did not carry.

This is where physical measurement becomes delicate. The temptation is to fill
every gap with a smooth story. A dashboard needle moves smoothly, so we imagine
speed itself has been recorded smoothly. A line graph joins dots, so we imagine
the line was observed. A simulation renders a continuous trajectory, so we
forget that the data entered at a finite rate. But the instrument did not give
us the continuum. It gave us marks. Everything between the marks is an
interpolation constrained by the marks and the calibration.

The closed question is:

\begin{center}
\emph{What can be said about the gap between recorded events without inventing
structure not in the ledger?}
\end{center}

The answer will not be ``nothing.'' That would be too weak. The gap has
boundary data. The marks on either side constrain any admissible continuation.
Nor will the answer be ``anything smooth.'' That would be too strong. Smoothness
can smuggle unmeasured structure into the record. The answer must lie between:
only continuations forced or permitted by the record, the calibration, and the
instrument's noise floor are admissible.

This chapter gives the gap a physical grammar. First, the speedometer shows
that missing information is a fact about resolution, not a mistake. Second,
bubble-chamber tracks show that interpolation is constrained by conservation
and apparatus geometry. Third, seismic data show how a minimum continuation can
be useful without pretending to be the event itself. Fourth, noise floors show
why some differences cannot enter the record. Finally, clinical-trial
metavariables show how an unresolved quantity can be carried without being
treated as resolved.

Chapter 1 gave us marks. Chapter 2 let instruments compare marks. Chapter 3
asks what comparison exposes: between marks there is absence, and absence has
structure.

\end{chapteressay}

% Phenomenon arcs use: 1/3 setup -> phenomenon box -> 2/3 structural interpretation.

\section{The Gap}
\label{se:the-gap}

The gap is the interval between admitted marks. It is not necessarily a temporal
interval, but in the speedometer it is easiest to see as time.

At highway speed, a Hall-effect sensor produces many pulses per second. If the
wheel and tooth count give roughly 470 pulses per second, the gap between
adjacent pulses is a little over two milliseconds. That is small for the driver,
but not zero. During that interval, the car continues to move. The tire deforms.
The engine torque changes. The road surface varies. The instrument carries no
new speed event until the next tooth crossing.

The display hides this discreteness. It may update smoothly, filter the pulse
stream, and show a stable number. The smooth display is a service to the driver,
not a literal copy of the record. The record is pulse, gap, pulse, gap. The
speed shown on the dashboard is an interpretation of that sequence.

This matters because events can fall inside the gap. A short skid, a vibration,
or a rapid correction may be physically real but absent from the speedometer's
ledger. A better instrument might catch it. A high-speed accelerometer might
see it. A tire sensor might see it. The speedometer does not. The gap is
instrument-local.

The right response is not to shame the speedometer. The instrument was never
asked to be a complete copy of the car. It was asked to produce an admissible
speed reading from finite marks. The gap is the price of that finitude.

This gives the chapter a first rule: absence from the ledger is not absence
from the world. It is absence from this instrument's carried record. The
difference is not academic. If a crash investigator reads only the speedometer
log, the investigator may miss a sub-interval event that an accelerometer
recorded. If an engineer reads only a sampled temperature trace, the engineer
may miss a fast transient. A physical record speaks only at the resolution and
through the carrier it has.

Yet the gap is not arbitrary. The mark before and the mark after restrict what
could have happened between them. A car cannot jump discontinuously across the
road in ordinary mechanics. A wheel cannot change rotation count without
passing through intermediate orientations. A sensor cannot receive a pulse that
its threshold did not admit. The gap is not blank paper; it is bounded paper.

The rest of this chapter is about that boundedness. What can the instrument say
without inventing? It can say what the marks force. It can say what the
calibration permits. It can say what the noise floor allows. It can carry a
placeholder for what remains unresolved. It cannot promote a convenient
interpolation into an observed fact.

\section{Interpolation as Physical Constraint}
\label{se:interpolation-as-physical-constraint}

Interpolation is often introduced as a mathematical convenience: draw a curve
between points. In measurement, interpolation is more serious. A drawn curve
can become a physical claim.

A bubble chamber makes the point. A charged particle passes through a
superheated liquid and leaves a trail of bubbles. The bubbles are not
continuous. They are discrete marks. A physicist looking at the photograph
draws a track through them. That track is not arbitrary. It is constrained by
the chamber geometry, the magnetic field, conservation of momentum, charge, and
the particle's possible interactions.

If the bubbles curve left in a magnetic field, the interpolation cannot curve
right merely because the drawn line would be pretty. If the bubbles are too far
apart, the uncertainty grows. If a decay occurs, the track may branch. The
physicist's line is an inference from marks under physical constraints.

The speedometer is simpler but not different in kind. Between pulses, the
instrument assumes a continuation. Over a short window, it may treat pulse rate
as stable enough to display a speed. During acceleration, filtering may lag.
During wheel slip, the continuation may be wrong. The interpolation is not a
free choice; it is part of the instrument's physical model.

This is why interpolation belongs inside measurement rather than after it. The
moment a display updates between admitted events, the instrument has already
chosen a continuation rule. The rule may be zero-order hold, moving average,
linear interpolation, spline smoothing, Kalman filtering, or something more
domain-specific. Each rule says what kind of unrecorded structure is allowed.

The constraint is negative as much as positive. A good interpolation does not
assert every possible detail. It refuses details the ledger cannot support. A
line between two points says: within the resolution of this record, and under
this model, these marks are connected this way. It should not say: this exact
path was observed.

That distinction will matter again in classical mechanics. A path selected by a
least-action principle is not a decorative curve between endpoints. It is the
admissible continuation forced by the record and the variational constraint.
But Chapter 6 cannot use that principle honestly unless Chapter 3 first names
the problem: between two records, there are many possible stories, and most of
them are inventions.

For now, the rule is modest. Interpolation is admissible only when it is
anchored at the recorded events, respects the calibration, and does not claim
more resolution than the instrument carried.

\section{The Minimum Admissible Continuation}
\label{se:the-minimum-admissible-continuation}

If many interpolations are possible, the instrument needs a discipline for
choosing among them. The discipline used throughout this project is minimality:
do not add structure that the record did not require.

Consider a seismologist with five pressure readings at five times. The readings
show a disturbance passing a station. A jagged line can connect the points. A
high-degree polynomial can pass through all of them while swinging wildly
between. A piecewise constant graph can hold each value until the next reading.
A cubic spline can join the points while minimizing unnecessary curvature. Each
candidate agrees with the marks. They do not make the same claim about the gap.

The minimum admissible continuation is not automatically the visually smoothest
curve. It is the continuation that introduces the least unsupported structure
while satisfying the constraints the record actually carries. If the apparatus
records only five values, a wildly oscillating curve between them is an
additional claim. It may be true if another instrument saw the oscillation, but
it is not justified by these five marks alone.

This is why spline methods are so natural in measurement. A spline can be read
as a disciplined refusal to invent extra bending. The endpoints and matching
conditions are honored; unnecessary curvature is suppressed. The curve is not
magic. It is a way of saying: here is the least-burden continuation compatible
with the recorded points and the smoothness conditions we are allowed to use.

The same logic underlies a speedometer's filtered display. A moving average or
low-pass filter suppresses unsupported high-frequency variation. It may hide
real rapid changes, so the rule must be chosen for the instrument's purpose. A
driver does not need a dashboard needle that jitters with every tooth. The
driver needs a stable estimate that does not invent precision.

Minimum continuation therefore has an ethical flavor. It is the discipline of
not pretending to know the gap better than the ledger allows. That discipline
does not make the continuation certain. It makes it admissible.

Later, this will become the heart of the selection problem. If a record admits
many continuations, which one is carried forward as the physical history? This
chapter does not yet answer that fully. It only prepares the ground: any answer
must be constrained by what is in the ledger and by a rule that prevents
unrecorded decoration.

\section{Noise and Signal}
\label{se:noise-and-signal}

The gap is not the only absence in measurement. Noise creates a second boundary:
the instrument may receive physical variation and still be unable to say which
part is signal.

A seismometer records ground motion continuously, but the record is never pure.
Ocean waves, wind, nearby traffic, thermal drift, electronic noise, and small
local vibrations all enter the apparatus. An earthquake signal must rise out of
that background. If the signal is below the noise floor, the instrument cannot
distinguish it. The event may have occurred physically, but the record cannot
carry it as a reliable mark.

The speedometer has a gentler version. The pulse stream may include jitter from
wheel vibration, small timing errors, electrical interference, or tire slip. The
ECU must decide which variation belongs to speed and which belongs to noise.
Filtering is not optional; without it the display would become unreadable. But
filtering is also a choice about the gap. It decides which changes are allowed
to survive.

The signal/noise distinction is therefore not a simple property of the world. It
is a property of the world as coupled to an instrument. A vibration that is noise
for the speedometer may be signal for a suspension diagnostic. A small ground
motion that is noise for one seismometer may be signal for a more sensitive
array. The boundary depends on the instrument's purpose, threshold, carrier, and
calibration.

This prevents a common overreach. One cannot say, from a flat record alone,
that nothing happened. One can say that nothing distinguishable from noise was
carried by this record. That statement is weaker, and again it is the useful
one. Physical inference is built from such weaker honest statements.

Noise also disciplines interpolation. A curve that follows every fluctuation may
be fitting noise rather than signal. A curve that suppresses too much may erase
real structure. The minimum admissible continuation must therefore respect the
noise floor. It should not chase variation the instrument cannot distinguish,
and it should not smooth away variation the instrument did carry.

In practice this is why measurement reports include resolution, uncertainty,
bandwidth, sampling rate, and filtering assumptions. Those quantities tell the
reader how to interpret the gap. They are not bureaucratic decorations. They are
part of the record's truth conditions.

\section{Metavariable in Measurement}
\label{se:metavariable-in-measurement}

Sometimes the gap is not a missing interval between two events. Sometimes it is
a named unknown carried by the record.

A clinical trial may report that a drug reduces systolic blood pressure by
$X$ mmHg, where $X$ is still under analysis. The trial has data. Patients were
enrolled, doses administered, pressures recorded, adverse events logged. But the
effect size has not yet been closed. The symbol $X$ is not empty. It is a
metavariable: a placeholder constrained by the record but not resolved into a
final value.

This is different from ignorance. Pure ignorance would say nothing. A
metavariable says: there is a value here, and the record constrains it, but the
analysis has not discharged the constraint. Acting as if $X=10$ before the
study closes is not using data; it is replacing a constrained placeholder with a
wish.

Measurement contains many such placeholders. A detector sees a candidate event
whose classification is pending. A calibration lab reports a provisional
standard awaiting final uncertainty analysis. A navigation system estimates a
position while satellite lock is incomplete. A weather model carries a
parameter whose value is being updated by new observations. In each case, the
instrument or study carries a slot that is not yet resolved.

The discipline is the same as for interpolation. The placeholder may be carried,
but it may not be treated as a completed mark. It can constrain later reasoning.
It can say which values are still possible and which values have been excluded.
It can guide additional measurement. It cannot close the ledger by fiat.

This is why the gap is productive. A missing value is not always a failure. It
can be a well-formed question. The art of measurement is to carry that question
without answering it prematurely.

The speedometer has a small version of this in its filtering window. Before
enough pulses arrive, the instrument may not yet have a stable rate. It can
carry a pending estimate. It can wait. It can update. If it displays too early,
it risks reporting a placeholder as a reading. If it waits too long, it becomes
unusable. The instrument's timing logic is therefore a policy for when a
metavariable has become a mark.

The chapter can now close. Between recorded events, the instrument may carry
constraints, estimates, noise floors, and placeholders. It may not invent
unrecorded structure and call it observation.

\begin{bridge}

The chapter's question was what can be said about the gap between recorded
events without inventing structure not in the ledger.

The answer is that the gap can be constrained but not filled freely. Recorded
marks provide boundary data. Calibration, conservation laws, apparatus geometry,
noise floors, and minimal-continuation rules restrict admissible interpolation.
Metavariables let the instrument carry unresolved values without pretending they
are resolved. The gap is therefore neither empty nor owned. It is a structured
absence.

Once gaps can be carried as structured absences, multiple records can begin to
interact. A speedometer, odometer, clock, accelerometer, radar trace, and GPS
track may each carry different marks and different gaps. Chapter 4 asks how
such records combine without collapsing into one undifferentiated record.

\end{bridge}

\begin{coda}{Security Camera Footage with Missing Frames}

The investigator sits in a small room behind the station office. On the monitor
is a corridor, a timestamp, and a person walking quickly toward an exit. The
camera records ten frames per second. The event in question lasted less than a
tenth of a second.

Frame 182 shows the person's hand just below the door handle. Frame 183 shows
the door already open. The question is whether the person used a key, forced
the latch, or found the door unsecured. The camera did not record the decisive
instant.

The missing frame becomes a metavariable.

The investigator cannot honestly draw any story at all into that gap. The
visible frames constrain the possibilities. The hand position in frame 182, the
door angle in frame 183, the amount of motion blur, the known stiffness of the
door closer, and the timing of the alarm contact all matter. If the alarm log
shows the latch released before the hand reached the handle, the story changes.
If the door frame shows damage, the story changes again. Each record constrains
the missing interval.

But the missing interval remains missing. The investigator may say that one
reconstruction is admissible and another is not. The investigator may say that
the available records do not distinguish between two possibilities. The
investigator may request another camera angle, a higher frame-rate recording, or
the electronic lock log. What the investigator may not do is promote a preferred
story into a recorded event.

The camera is not defective because it has gaps. Its frame rate is the
condition under which it can carry a record at all. A faster camera would have
different gaps. A slower camera would have larger ones. The forensic question is
not how to erase the gap, but how to respect it.

Noise appears too. Compression artifacts blur the hand. Fluorescent lights
flicker. Shadows move across the door. A reflection in the glass looks like a
second person until another angle shows it is only the corridor behind the
camera. The investigator learns to separate signal from noise by comparing
records. The visible image alone is not sovereign.

The security footage performs the chapter's structure. The frames are admitted
marks. The missing instant is the gap. The reconstruction is an interpolation
constrained by boundary data. The unresolved action at the handle is a
metavariable. The alarm log, door damage, and second camera are additional
records that may reduce the gap but do not retroactively make the first camera
continuous.

At the end of the report, the investigator writes carefully: the footage is
consistent with a key being used, but does not prove it; the forced-entry
hypothesis is inconsistent with the lack of frame damage; the lock log is needed
to close the question.

That is not weakness. That is measurement telling the truth about its gaps.

\end{coda}
